\documentclass[a4paper, serif, 10pt]{moderncv}

%% moderncv
% color and style
\moderncvcolor{black}
\moderncvstyle{banking}

%% page layout/margins
\usepackage[top=2.5cm, bottom=2.5cm, left=1.5cm, right=1.5cm]{geometry}

\nopagenumbers{}

%% Babel config for English language
\usepackage[english]{babel}

%% fontspec
\usepackage{fontspec}

\IfFontExistsTF{Linux Libertine}{
  \setmainfont[Ligatures={TeX, Common}, Numbers=OldStyle]{Linux Libertine}
}{}
\IfFontExistsTF{Linux Libertine O}{
  \setmainfont[Ligatures={TeX, Common}, Numbers=OldStyle]{Linux Libertine O}
}{}

%% CTeX
% CJK support, used to show CJK characters in the resume
%
% - heading=false: disable ctex heading
% - punct=kaiming: use kaiming punctuations styles for CJK
% - scheme=plain: use plain scheme, do not override \normalsize font size
% - space=auto: space settings for CJK characters
%
% CJK fonts are resolved per locale language, e.g. Noto Serif/Sans CJK JP
% for Japanese resumes, with fallback to the next region if not installed
%
% ref:
% - http://ctan.mirrorcatalogs.com/language/chinese/ctex/ctex.pdf
\usepackage[UTF8, heading=false, punct=kaiming, scheme=plain, space=auto]{ctex}

\IfFontExistsTF{Noto Serif CJK SC}{
  \setCJKmainfont{Noto Serif CJK SC}
}{}
\IfFontExistsTF{Noto Sans CJK SC}{
  \setCJKsansfont{Noto Sans CJK SC}
}{}

%% line spacing
\usepackage{setspace}
\setstretch{1.125}

%% URL styling - use normal text instead of monospace
\urlstyle{same}

%% Patch \httplink and \httpslink to handle full URLs
% The original moderncv macros always prepend http:// or https:// to URLs.
% This causes issues when the URL already has a protocol (e.g., https://example.com)
% resulting in malformed URLs like https://https://example.com.
% This patch checks if the URL already contains "://" and uses it directly if so.
% Note: We use \str_set:Nx (x-expansion) to expand macros like \@homepage first.
\ExplSyntaxOn
\str_new:N \g_yamlresume_url_str

\RenewDocumentCommand{\httpslink}{O{}m}{
  \str_set:Nx \g_yamlresume_url_str {#2}
  \str_if_in:NnTF \g_yamlresume_url_str {://}
    { \url{#2} }
    { \url{https://#2} }
}

\RenewDocumentCommand{\httplink}{O{}m}{
  \str_set:Nx \g_yamlresume_url_str {#2}
  \str_if_in:NnTF \g_yamlresume_url_str {://}
    { \url{#2} }
    { \url{http://#2} }
}
\ExplSyntaxOff

\name{Giulia Bianchi}{}
\title{Responsabile Risorse Umane e Sviluppo del Talento}
\phone[mobile]{+39 02 555 0182}
\email{giulia.bianchi@giuliabianchi-hr.it}
\homepage{https://www.giuliabianchi-hr.it}

\address{Via Alessandro Manzoni 12, Milano, Lombardia, Italy, 20121}{}{}

\extrainfo{{\small \faLinkedin}\ \href{https://www.linkedin.com/in/giulia-bianchi-hr}{@giulia-bianchi-hr}}

\begin{document}

\maketitle

\section{Basics}

\cvline{}{Responsabile Risorse Umane con oltre dieci anni di esperienza nella guida di gruppi di lavoro e nella definizione di strategie orientate alle persone.
%
\begin{itemize}
\item Gestione completa del ciclo del talento: selezione, inserimento, formazione, valutazione delle prestazioni e sviluppo delle carriere.
\item Esperienza consolidata nelle relazioni sindacali, nella contrattazione collettiva e nella conformità normativa.
\item Capacità di tradurre gli obiettivi aziendali in politiche e strumenti concreti per le Risorse Umane.
\item Orientamento ai risultati, alla qualità delle relazioni interne e al benessere organizzativo.
\end{itemize}}
\section{Education}

\cventry{Sep 2009–Jul 2011}
        {Master, Economia e Gestione delle Risorse Umane, Score: 110/110 con lode}
        {Università Commerciale Luigi Bocconi}
        {\url{https://www.unibocconi.it}}
        {}
        {\begin{itemize}
\item Formazione avanzata nella gestione strategica delle Risorse Umane e nelle relazioni industriali.
\item Partecipazione a laboratori e progetti con aziende del territorio su selezione, valutazione e sviluppo del personale.
\item Tesi sperimentale sui modelli di valutazione delle prestazioni nelle imprese di servizi.
\end{itemize}
\textbf{Courses}: Gestione Strategica delle Risorse Umane, Diritto del Lavoro e Relazioni Industriali, Comportamento Organizzativo, Economia del Personale, Sistemi di Valutazione delle Prestazioni, Statistica per le Scienze Sociali}

\cventry{Sep 2006–Jul 2009}
        {Bachelor, Scienze dell'Organizzazione e del Lavoro, Score: 108/110}
        {Università degli Studi di Milano}
        {\url{https://www.unimi.it}}
        {}
        {\begin{itemize}
\item Acquisizione di solide basi in sociologia del lavoro, psicologia delle organizzazioni e diritto sindacale.
\item Sviluppo di competenze analitiche attraverso ricerche sul campo e analisi dei dati organizzativi.
\end{itemize}
\textbf{Courses}: Sociologia del Lavoro, Psicologia delle Organizzazioni, Diritto Sindacale, Economia Aziendale, Metodologia della Ricerca Sociale}
\section{Work}

\cventry{Jan 2020–Present}
        {Responsabile Risorse Umane}
        {Lombardia Energia S.p.A.}
        {\url{https://www.lombardiaenergia.it}}
        {}
        {\begin{itemize}
\item Guida di un gruppo di 6 professionisti delle Risorse Umane, con responsabilità su selezione, formazione, sviluppo e amministrazione del personale.
\item Definizione e attuazione della strategia delle Risorse Umane in linea con gli obiettivi aziendali e con il piano industriale.
\item Gestione delle relazioni sindacali e della contrattazione collettiva, con riduzione del contenzioso del 30\% in tre anni.
\item Introduzione di un modello di valutazione delle prestazioni e dei percorsi di carriera, aumentando la mobilità interna del 25\%.
\item Progettazione di programmi di benessere aziendale e di diversità e inclusione, con un incremento del coinvolgimento dei dipendenti di 18 punti percentuali.
\item Monitoraggio degli indicatori chiave di prestazione delle Risorse Umane e predisposizione della reportistica per la direzione.
\end{itemize}
\textbf{Keywords}: Strategia delle Risorse Umane, Relazioni sindacali, Sviluppo del talento, Benessere aziendale, Diversità e inclusione, Budget del personale}

\cventry{Mar 2016–Dec 2019}
        {Referente Risorse Umane per le Funzioni Aziendali}
        {TechNova Italia S.r.l.}
        {\url{https://www.technovaitalia.it}}
        {}
        {\begin{itemize}
\item Partner di riferimento per le funzioni commerciale e operativa nella gestione delle tematiche delle Risorse Umane.
\item Gestione completa dei processi di selezione per profili tecnici e manageriali.
\item Supporto ai responsabili di funzione nella valutazione delle prestazioni, nella definizione degli obiettivi e nei piani di sviluppo individuali.
\item Progettazione ed erogazione di percorsi formativi, con oltre 1.200 ore di formazione erogate in tre anni.
\item Implementazione di un nuovo sistema informativo del personale e digitalizzazione dei processi di inserimento e valutazione.
\item Partecipazione a progetti di riorganizzazione aziendale e gestione del cambiamento.
\end{itemize}
\textbf{Keywords}: Selezione del personale, Valutazione delle prestazioni, Formazione, Sistemi del personale, Gestione del cambiamento}

\cventry{Sep 2011–Feb 2016}
        {Specialista Sviluppo Organizzativo}
        {Gruppo Editoriale Aurora S.p.A.}
        {\url{https://www.gruppoeditorialeaurora.it}}
        {}
        {\begin{itemize}
\item Analisi dei fabbisogni organizzativi e progettazione di percorsi di sviluppo e carriera per circa 400 dipendenti.
\item Realizzazione di indagini sul clima aziendale e definizione di piani di azione per migliorare il benessere organizzativo.
\item Mappatura delle competenze e supporto alla revisione del modello organizzativo.
\item Coordinamento di attività formative e gestione dei rapporti con enti di formazione esterni.
\item Contributo alla redazione di policy interne in materia di gestione delle Risorse Umane.
\end{itemize}
\textbf{Keywords}: Analisi dei fabbisogni, Percorsi di carriera, Clima aziendale, Formazione, Mappatura delle competenze}
\section{Languages}

\cvline{Italian}{Native or Bilingual Proficiency \hfill \textbf{Keywords}: Madrelingua}
\cvline{English}{Full Professional Proficiency \hfill \textbf{Keywords}: Certificazione C1, Inglese commerciale}
\cvline{French}{Limited Working Proficiency \hfill \textbf{Keywords}: B1, Conversazione professionale}
\section{Skills}

\cvline{Gestione delle Risorse Umane}{Expert \hfill \textbf{Keywords}: Selezione del personale, Valutazione delle prestazioni, Sviluppo del talento, Pianificazione della successione, Benessere aziendale}
\cvline{Diritto del Lavoro e Relazioni Sindacali}{Advanced \hfill \textbf{Keywords}: Contrattazione collettiva, Relazioni sindacali, Gestione delle controversie, Conformità normativa}
\cvline{Sistemi Informativi del Personale}{Advanced \hfill \textbf{Keywords}: SAP SuccessFactors, Zucchetti, Excel avanzato, Power BI, Sistemi del personale}
\cvline{Sviluppo Organizzativo}{Advanced \hfill \textbf{Keywords}: Analisi dei fabbisogni, Mappatura delle competenze, Gestione del cambiamento, Clima aziendale, Diversità e inclusione}
\cvline{Competenze Trasversali}{Expert \hfill \textbf{Keywords}: Comunicazione efficace, Negoziazione, Capacità di guida, Risoluzione dei problemi, Pensiero analitico}
\section{Awards}

\cventry{Nov 2022}
        {Associazione Italiana per la Direzione del Personale}
        {Premio Eccellenza Risorse Umane}
        {}
        {}
        {Riconoscimento assegnato a professionisti che si sono distinti per innovazione nella gestione delle Risorse Umane e per l'impatto misurabile sui risultati aziendali.}
\section{Certificates}

\cventry{Mar 2021}
        {Istituto Italiano di Certificazione delle Risorse Umane}
        {Certificazione Professionista Senior Risorse Umane}
        {\url{https://www.iicru.it/certificazioni}}
        {}
        {}

\cventry{Jun 2018}
        {Ordine dei Consulenti del Lavoro di Milano}
        {Esperto in Diritto del Lavoro}
        {\url{https://www.consulentidellavoro.mi.it}}
        {}
        {}
\section{Publications}

\cventry{Apr 2023}
        {Gestione del talento nell'era del lavoro ibrido}
        {Rivista Italiana di Direzione del Personale}
        {\url{https://www.rivistaitalianadirigenzapersonale.it/articoli/talento-lavoro-ibrido}}
        {}
        {\begin{itemize}
\item Analisi delle principali tendenze del lavoro ibrido e delle implicazioni per la gestione del talento.
\item Proposta
\end{itemize}}
\end{document}